\title{Direct Preference Optimization: Your Language Model is Secretly a Reward Model}

\begin{abstract}
Although large unsupervised language models (LMs) acquire extensive world knowledge and display some reasoning capabilities, fine-grained management of their outputs remains challenging due to the fully unsupervised character of their initial training. Current approaches aiming to enhance model controllability typically involve collecting human preference data on generation quality and adjusting the unsupervised LM via fine-tuning to match these human judgements. The predominant strategy, reinforcement learning from human feedback (RLHF), entails a multi-stage and frequently unstable pipeline: first, learning a reward model that encodes human evaluations, and subsequently optimizing the base LM through reinforcement learning to maximize the estimated reward while maintaining fidelity to the original parameters. In this work, we propose a novel reward model parameterization for RLHF, which permits explicit derivation of the associated optimal policy, thereby transforming the canonical RLHF challenge into a tractable classification loss optimization. Our proposed method, termed \textit{Direct Preference Optimization} (DPO), is robust, efficient, and computationally economical, obviating the need for LM sampling during training or extensive hyperparameter searches. Through comprehensive empirical evaluation, we demonstrate that DPO is capable of aligning LMs to human preferences at least as effectively as alternative techniques. Importantly, DPO surpasses RLHF based on PPO for sentiment control tasks and either matches or outperforms it in summarization and single-turn dialogue, all while offering significant advantages in simplicity of training and implementation.
\end{abstract}

\section{Introduction}
Large-scale unsupervised language models (LMs), when trained on vast amounts of data, often demonstrate unexpected abilities~\citep{chowdhery2022palm, brown2020language, touvron2023llama,bubeck2023sparks}. Yet, because these models are exposed to human-generated data characterized by a diverse spectrum of intentions, expertise, and objectives, they inevitably assimilate both desirable and undesirable patterns. For instance, although it is beneficial for an AI coding assistant to \textit{recognize} frequently occurring programming errors to help rectify them, we would prefer that the model preferentially emulate the relatively rare but expert-level code present in its training distribution. Likewise, a language model should be \textit{informed} about prevalent misconceptions—such as those held by half the population—without echoing these inaccuracies in half of its responses to related queries. Put differently, distilling the LM’s \emph{target behaviors and answers} from the extensive pool of its \textit{learned knowledge and skills} is essential for constructing AI agents that are reliable, high-performing, and easy to direct~\citep{ouyang2022training}. Conventional approaches primarily attempt to align LMs with human values through reinforcement learning (RL). Here, we demonstrate that the standard RL-driven objective can be exactly reformulated as a binary cross-entropy objective, thereby streamlining the preference alignment process.

Typically, preference alignment in language models proceeds by introducing training signals derived from curated human preference datasets, specifying the desired behavioral tendencies such as safety and helpfulness. This stage succeeds the initial unsupervised pre-training, which exposes the model to a broad corpus of text. While an intuitive solution is to perform supervised fine-tuning on exemplary human-generated outputs, current state-of-the-art methods predominantly rely on reinforcement learning from human or AI-derived feedback (RLHF/RLAIF; \citep{christiano2017deep,bai2022constitutional}). In RLHF, a reward model is first fitted to human preference data, which then guides an RL algorithm to adapt the language model’s policy, maximizing the generation of high-reward outputs while constraining deviations from the original policy. Although RLHF has enabled the emergence of powerful conversational and programming capabilities in language models, this methodology introduces considerable complexity relative to standard supervised learning, as it demands training multiple LMs and incorporates on-policy sampling, thus resulting in a substantial computational burden.

This work introduces a method for training language models to align with human preferences without the need for explicit reward modeling or reinforcement learning frameworks. We present \textit{Direct Preference Optimization} (DPO), an approach that implicitly targets the same optimization goal as traditional RLHF methods (maximizing rewards while enforcing a KL-divergence regularization), yet offers a much more straightforward and easily trainable formulation. Conceptually, the DPO method updates the model by amplifying the log probability of preferred outputs relative to less desirable ones, while leveraging a dynamically computed, instance-specific importance weight to guard against the model collapse that can arise from naively applying a probability ratio objective. DPO, like prior techniques, is grounded in a formal preference model (for example, the Bradley-Terry model; \cite{bradley1952rankanalysis}), which serves to quantify the agreement between a candidate reward function and observed preference judgments. However, rather than using this preference model to supervise reward model training followed by subsequent policy optimization, DPO employs a variable substitution to express the preference loss in terms of the policy itself. This formulation enables DPO to optimize the policy directly using a basic binary cross-entropy loss over datasets of human response preferences, yielding an optimal policy corresponding to an implicit reward inferred from these preferences.

The central contribution of this work is the introduction of Direct Preference Optimization (DPO), a reinforcement learning-free and straightforward algorithm for preference-based language model training. Through empirical evaluation, we demonstrate that DPO performs on par with, or better than, leading methods—such as RLHF approaches based on PPO—in a variety of tasks including sentiment adjustment, text summarization, and dialogue generation, even when scaling to language models with up to 6B parameters.

Language models trained via self-supervised learning at increasing scales have demonstrated the capacity to address certain tasks in a zero-shot setting \citep{radford2019language}, or through few-shot prompting strategies \citep{gpt3,megatron,chowdhery2022palm}. Nevertheless, further performance enhancements on downstream applications and stronger alignment with user objectives can be achieved by fine-tuning on corpora composed of instructions and human-generated completions \citep{mishra-etal-2022-cross,sanh2022multitask,chung2022scaling,thoppilan2022lamda}. This process, known as `instruction-tuning,' not only improves the language models' capability to follow new, unseen instructions but also broadly enhances their practical usability \citep{chung2022scaling}. While instruction tuning has proven effective, acquiring \textit{relative} human assessments of response quality is generally less resource-intensive than collecting expert-level demonstrations. Consequently, subsequent research has explored fine-tuning large language models using datasets annotated with human preferences, resulting in gains across tasks including translation \citep{kreutzer-etal-2018-reliability}, summarization \citep{stiennon2022learning,ziegler2020finetuning}, story-telling \citep{ziegler2020finetuning}, and instruction following \citep{ouyang2022training,ramamurthy2023is}. Approaches in this area typically begin by training a neural reward model to match the distribution of observed preferences, frequently employing the Bradley-Terry preference model \citep{bradley1952rankanalysis}. The language model is then fine-tuned through reinforcement learning techniques—such as REINFORCE \citep{williams1992reinforce}, proximal policy optimization (PPO; \cite{schulman2017proximal}), or related algorithms \citep{ramamurthy2023is}—to maximize this learned reward. A related research trajectory utilizes instruction-following LLMs refined with human feedback to generate additional synthetic data reflecting preferences for particular qualities like safety or harmlessness \citep{bai2022constitutional}, relying on limited human oversight in the form of rubric-based textual evaluations. These methodologies sit at the intersection of two main research threads: one focused on optimizing language models using reinforcement learning for various goals~\citep{Ranzato2015SequenceLT,paulus2018a,wu2018learning}, and another centered on methods for learning directly from human preference signals \citep{christiano2017deep,kupcsik2018learning}. Notwithstanding the advantages of utilizing comparative human preferences, applying reinforcement learning-based fine-tuning to large language models presents considerable practical difficulties; the current work advances a principled method for preference optimization that circumvents the need for RL.

Beyond the domain of language, the study of policy learning from preference feedback has been explored within both bandit and reinforcement learning frameworks, resulting in a variety of proposed methods. When action preferences or rankings, rather than explicit rewards, are utilized in contextual bandit problems, the formulation is referred to as a contextual dueling bandit (CDB; \cite{yue2012karmed,dudik2015contextual}). In scenarios lacking absolute reward information, the theoretical foundations of CDBs employ the concept of a \textit{von Neumann winner}—a policy whose expected probability of outperforming any other policy meets or exceeds 50\% \citep{dudik2015contextual}. Notably, CDB frameworks assume that preference signals are provided sequentially in an online fashion, whereas in human preference learning, the common setting involves learning from a static collection of offline preference-labeled action pairs \citep{yan2022human}. Along similar lines, \textit{preference-based RL} (PbRL) makes use of binary preferences elicited from an unknown underlying `scoring' function, in place of direct reward signals \citep{BusaFekete2014,ruiz2023dueling}. There exist numerous algorithms for PbRL, many of which are capable of leveraging off-policy preference information; however, most require an explicit estimation of the underlying scoring or reward model before proceeding with policy optimization \citep{jain2013learning,BusaFekete2014,christiano2017deep,sadigh2017active,kupcsik2018learning}. In contrast, we introduce a unified policy learning framework that bypasses intermediate reward modeling and directly optimizes the policy according to observed preferences.

\section{Preliminaries}\label{section:prelims}

We provide an overview of the RLHF procedure as outlined by \citeauthor{ziegler2020finetuning} and subsequently by works such as \citep{stiennon2022learning, bai2022training, ouyang2022training}. The standard RLHF workflow generally consists of three main components: 1) supervised fine-tuning (SFT); 2) preference collection and reward model training; and 3) policy optimization via RL.

\textbf{SFT}: The RLHF process customarily starts by adapting a pre-trained language model (LM) to the specific task (such as dialogue or summarization) through supervised learning on curated datasets, resulting in an initial model $\pi^\text{SFT}$.

\textbf{Reward Modelling Phase}: During this stage, prompts $x$ are given to the SFT model, which produces response pairs $(y_1, y_2)\sim \pi^\text{SFT}(y \mid x)$. Human evaluators then assess these pairs by indicating a preferred response, captured as $y_w\succ y_l \mid x$, with $y_w$ signifying the chosen completion and $y_l$ the rejected alternative among $(y_1, y_2)$. These recorded preferences are assumed to reflect judgments arising from a hidden reward function $r^*(y, x)$, to which we lack direct access. A variety of probabilistic models exist for mapping preferences to reward, with the Bradley-Terry (BT) model \cite{bradley1952rankanalysis} frequently employed (alternatively, more general Plackett-Luce models \citep{plackett1975analysis, luce2012individual} may be appropriate when multiple ranked completions are available). Under the BT model, the human preference probability distribution $p^*$ can be expressed as:

Given a static comparison dataset $\mathcal{D} = \bigl\{x^{(i)}, y_w^{(i)}, y_l^{(i)}\bigr\}_{i=1}^N$ drawn from $p^*$, we can specify a parameterized reward function $r_{\phi}(x, y)$ and fit its parameters using maximum likelihood estimation. If the problem is framed as binary classification, the negative log-likelihood loss is:

\begin{equation}\label{eq:reward_model}
    \mathcal{L}_R(r_{\phi}, \mathcal{D}) = -\mathbb{E}_{(x, y_w, y_l)\sim \mathcal{D}}\bigl[\log \sigma(r_{\phi}(x, y_w)- r_{\phi}(x, y_l))\bigr]
\end{equation}

with $\sigma$ representing the logistic sigmoid. When applying this to LMs, it is standard to initialize $r_{\phi}(x, y)$ from an SFT model $\pi^\text{SFT}(y \mid x)$, followed by a linear projection atop the last transformer block, outputting a scalar reward estimate \cite{ziegler2020finetuning}. To encourage stability in the reward function, previous research typically enforces a normalization constraint, setting $\mathbb{E}_{x,y\sim \mathcal{D}}\left[r_\phi(x, y)\right] = 0$ across $x$.

\textbf{RL Fine-Tuning Phase}: During reinforcement learning, the estimated reward is used to guide updates to the language model policy. Following established work~\citep{jaques2017sequence, jaques2020human}, optimization is carried out as follows:

\begin{equation}\label{eq:RL}
\max_{\pi_{\theta}}  \mathbb{E}_{x\sim \mathcal{D}, y\sim \pi_{\theta}(y \mid x)}\bigl[r_{\phi}(x, y)\bigr] - \beta\mathbb{D}_{\textrm{KL}}\bigl[\pi_{\theta}(y\mid x)\mid \mid \pi_\text{ref}(y\mid x)\bigr],
\end{equation}

where $\beta$ determines the permissible divergence from a fixed reference policy $\pi_\text{ref}$, usually taken to be the pre-trained SFT policy $\pi^\text{SFT}$. Practically, $\pi_\theta$ is initialized from $\pi^\text{SFT}$ as well. This regularization is critical: it constrains the updated policy to stay near the data regime where the reward model is reliable, maintains generation variety, and avoids mode collapse on a narrow subset of high-reward outputs. Given that generation is discrete, direct gradient methods are not applicable; thus, policy optimization is generally achieved via reinforcement learning. The common practice~\citep{ziegler2020finetuning, stiennon2022learning, bai2022training, ouyang2022training} is to define the reward for policy optimization as $r(x, y) = r_{\phi}(x, y) -\beta (\log \pi_{\theta}(y\mid x) - \log \pi_\text{ref}(y\mid x))$ and to maximize it using PPO \cite{schulman2017proximal}.

\section{Direct Preference Optimization}\label{sec:DPO}

Recognizing the complexity of employing reinforcement learning on large-scale language model fine-tuning, our objective is to develop a direct and efficient method for policy optimization from preference data. Unlike RLHF pipelines that separately train a reward model and then employ RL, our methodology is based on a specific form of reward model parameterization that makes it possible to analytically derive the associated optimal policy, circumventing the need for a full RL process. We detail below how an analytical relationship between reward models and their corresponding optimal policies allows us to re-express a loss on rewards in terms of policy parameters. This variable substitution approach negates the necessity of explicitly training a reward model, yet still ensures adherence to standard preference models such as the Bradley-Terry framework. Consequently, the policy model simultaneously acts as both the generator and the implicit reward estimator.

\textbf{DPO Objective Derivation.} We begin from the same reinforcement learning (RL) formulation as existing studies, specified in Eq.~\ref{eq:RL}, with an arbitrary reward function $r$. As established in prior literature~\citep{peters2007reinforcement, peng2019advantage, korbak2022reinforcement, go2023aligning}, it is not difficult to show that the solution maximizing expected reward subject to a KL-divergence constraint on policy, as in Eq.~\ref{eq:RL}, is given by:

\begin{equation}\label{eq:op_policy}
    \pi_r(y\mid x) = \frac{1}{Z(x)}\pi_\text{ref}(y\mid x)\exp\left(\frac{1}{\beta}r(x, y)\right),
\end{equation}

where $Z(x) =\sum_{y}\pi_\text{ref}(y\mid x)\exp\left(\frac{1}{\beta}r(x, y)\right)$ serves as the normalization constant. A full derivation is provided in Appendix \ref{app:derivation1}. Although estimating $r^*$ via MLE to yield $r_{\phi}$, the computation of $Z(x)$ remains computationally prohibitive \citep{korbak2022reinforcement, go2023aligning}, limiting the practicality of directly using this expression. However, we can rewrite Eq.~\ref{eq:op_policy} to isolate the reward as a function of the optimal policy $\pi_r$, the baseline policy $\pi_\text{ref}$, and the still intractable $Z(x)$. By taking the logarithm of both sides and rearranging, we find:

\begin{equation}\label{eq:main_eq}
    r(x,y) =\beta \log \frac{\pi_r(y\mid x)}{\pi_\text{ref}(y\mid x)} + \beta \log Z(x).
\end{equation}

Applying this transformation to the actual reward function $r^*$ and its associated optimal policy $\pi^*$ is straightforward. Notably, the Bradley-Terry framework relies solely on the reward difference between two outputs, i.e., ${p^*(y_1 \succ y_2 \mid x) = \sigma(r^*(x, y_1) - r^*(x, y_2))}$. Upon substituting the reformulated reward from Eq.~\ref{eq:main_eq} into the Bradley-Terry preference formulation (Eq.~\ref{eq:bradley-terry}), the partition function $Z(x)$ is eliminated, making it possible to express the human preference probability exclusively in terms of $\pi^*$ and $\pi_\text{ref}$. Thus, the Bradley-Terry model yields that the optimal RLHF policy $\pi^*$ satisfies:

\begin{equation}\label{eq:objective}
    p^*(y_1\succ y_2 \mid x)=\frac{1}{1 + \exp\left(\beta \log \frac{\pi^*(y_2\mid x)}{\pi_\text{ref}(y_2\mid x)} - \beta \log \frac{\pi^*(y_1\mid x)}{\pi_\text{ref}(y_1\mid x)}\right)}
\end{equation}

A complete derivation is in Appendix~\ref{app:derivation2}. Although Eq.~\ref{eq:objective} focuses on the Bradley-Terry formulation, a comparable reparameterization applies to the broader Plackett-Luce family~\citep{plackett1975analysis, luce2012individual}; see Appendix~\ref{app:plackett_luce_models} for those derivations.

With human preference probabilities now parameterized via the optimal policy (rather than the reward), we are able to propose a maximum likelihood estimation procedure for any parameterized policy $\pi_\theta$. In analogy to the approach for fitting reward models (cf. Eq.~\ref{eq:reward_model}), we define our policy training objective as:

\begin{equation}\label{eq:optimum_model}
    \mathcal{L}_\text{DPO}(\pi_{\theta}; \pi_\text{ref}) = -\mathbb{E}_{(x, y_w, y_l)\sim \mathcal{D}}\left[\log \sigma \left(\beta \log \frac{\pi_{\theta}(y_w\mid x)}{\pi_\text{ref}(y_w\mid x)} - \beta \log \frac{\pi_{\theta}(y_l\mid x)}{\pi_\text{ref}(

By employing this approach, we learn an implicit reward through a different parameterization, for which the corresponding optimal policy directly corresponds to $\pi_\theta$. Additionally, because our method aligns with fitting a reparametrized Bradley-Terry model, it possesses desirable theoretical guarantees, such as consistency given certain assumptions about the preference data distribution \cite{bong2022generalized}. We further elaborate on the theoretical characteristics of DPO and compare it to related literature in Section~\ref{sec:theory}.

\textbf{What mechanism underlies the DPO update?} To gain insight into the operation of DPO, one may examine the gradient of its loss function $\mathcal{L}_\text{DPO}$. The derivative with respect to the model parameters $\theta$ is given by:

\begin{multline*}\label{eq:gradient}
    \nabla_\theta \mathcal{L}_\text{DPO}(\pi_\theta;\pi_\text{ref}) = \\ -\beta\mathbb{E}_{(x, y_w, y_l) \sim \mathcal{D}} \bigg[\underbrace{\sigma(\hat{r}_\theta(x, y_l) - \hat{r}_\theta (x, y_w))}_\text{greater influence when reward estimates are inaccurate}\bigg[\underbrace{\nabla_\theta\log \pi(y_w \mid x)}_\text{push up probability of $y_w$} - \underbrace{\nabla_\theta\log\pi(y_l \mid x)}_\text{push down probability of $y_l$}\bigg]\bigg],
\end{multline*}

Here, $\hat{r}_\theta(x, y) = \beta \log \frac{\pi_\theta(y \mid x)}{\pi_\text{ref}(y \mid x)}$ denotes the reward implicitly determined by the parameterized language model $\pi_\theta$ in reference to $\pi_\text{ref}$ (see Section~\ref{sec:theory} for further details). The essential effect of the gradient of $\mathcal{L}_\text{DPO}$ is to promote the probability of selecting the favored outputs $y_w$, while discouraging the selection of less preferred outputs $y_l$. Crucially, each example is weighted in accordance with how much the implicit reward model $\hat{r}_\theta$ overestimates the reward of dispreferred outputs compared to preferred ones, scaled by $\beta$; in other words, the weighting reflects the degree to which the reward model fails to order completions properly, factoring in the impact of the KL regularization. Experimental findings highlight the significance of this weighting scheme: omitting it (i.e., using a simplistic variant without the weighting term) leads to language model degradation (see Appendix Table~\ref{tab:unlikelihood_generations}).

\textbf{DPO Overview.} The typical DPO workflow consists of the following steps: First, for each prompt $x$, generate completion samples $y_1, y_2 \sim \pi_\text{ref}(\cdot \mid x)$, then annotate these pairs based on human preference judgments, thus forming the offline preference dataset $\mathcal{D} = \{x^{(i)}, y_w^{(i)}, y_l^{(i)}\}_{i=1}^N$. Second, the language model $\pi_\theta$ is trained to minimize the loss $\mathcal{L}_\text{DPO}$ with respect to the chosen $\pi_\text{ref}$, the dataset $\mathcal{D}$, and the desired temperature parameter $\beta$.

In practical scenarios, it is often preferable to leverage existing, publicly available preference datasets instead of collecting new samples and human annotations. Because these preference datasets are typically derived using $\pi^\text{SFT}$, we adopt $\pi_\text{ref} = \pi^\text{SFT}$ wherever this distribution is accessible. If $\pi^\text{SFT}$ cannot be used, we instead construct $\pi_\text{ref}$ by maximizing the likelihood of the preferred responses, i.e., ${\pi_\text{ref} = \argmax_{\pi}\mathbb{E}_{x, y_w \sim \mathcal{D}}\left[\log \pi(y_w \mid x)\right]}$. This approach serves to minimize the distribution mismatch between the actual (but unknown) reference distribution and the $\pi_\text{ref}$ employed by DPO. Additional details on implementation specifics and the selection of hyperparameters are provided in Appendix~\ref{app:implementation}.

\section{Theoretical Analysis of DPO} This section provides a deeper examination of the DPO technique, elaborates on the theoretical foundations, and clarifies the benefits of DPO in relation to known limitations of actor-critic approaches used in RLHF settings, such as PPO~\cite{schulman2017proximal}.

\label{sec:theory}

\subsection{Your Language Model Is Secretly a Reward Model} DPO circumvents the need to explicitly learn a reward function and avoids subsequent policy optimization through reinforcement learning, by directly employing a single maximum likelihood formulation. Specifically, the loss function given in Eq.~\ref{eq:main_eq} aligns with the Bradley-Terry model, where the reward is parameterized as $r^*(x, y) = \beta \log\frac{\pi^*_\theta(y \mid x)}{\pi_\text{ref}(y \mid x)}$, and we train the parameterized model $\pi_{\theta}$ in a manner analogous to optimizing a reward model as in Eq.~\ref{eq:reward_model}, after an appropriate change of variables. The subsequent theoretical development establishes the validity of this reparameterization, demonstrates that it does not impose limitations on the representational capacity of learned reward models, and proves that it enables exact recovery of the optimal policy. We start by introducing an equivalence relation among reward functions.

\begin{definition}
We define two reward functions $r(x, y)$ and $r'(x, y)$ as equivalent if and only if ${r(x, y)-r'(x, y) = f(x)}$ for some function $f$.
\end{definition}

It is straightforward to verify that this relationship forms an equivalence relation, which thus divides the set of all reward functions into equivalence classes. From this, we can state the following lemmas:

\begin{lemma}\label{lemma:same_prefrence} Within the Plackett-Luce framework, including its special case Bradley-Terry, two reward functions belonging to the same equivalence class induce identical preference distributions.
\end{lemma}

\begin{lemma}\label{lemma:same_policy}
    Any two reward functions from a shared equivalence class lead to the same optimal policy in the setting of constrained RL.
\end{lemma}

The arguments establishing these results are routine, and thus their proofs are postponed to Appendix \ref{app:lemma1}. The first lemma points to a recognized under-specification problem in the Plackett-Luce class of models \cite{plackett1975analysis}. Owing to this ambiguity, it is customary to introduce supplementary identifiability conditions in order to derive well-defined MLE estimates from Eq. \ref{eq:reward_model} \cite{bong2022generalized}. The second lemma asserts that the choice of reward function within a single class does not affect the optimal policy; accordingly, our focus reduces to identifying any representative reward function within the optimal equivalence class. The following theorem, which we demonstrate in Appendix~\ref{app:thm1}, makes this precise:

\begin{theorem}\label{thm:main}
    Given certain mild regularity conditions, any class of rewards consistent with Plackett-Luce (and specifically Bradley-Terry) models can be parameterized as ${r(x, y) = \beta \log \frac{\pi(y\mid x)}{\pi_\text{ref}(y\mid x)}}$ for some model $\pi(y\mid x)$ and chosen reference $\pi_\text{ref}(y \mid x)$.
\end{theorem}

\begin{sproof}
    Let us consider an arbitrary reward function $r(x, y)$ that defines an associated optimal model $\pi_r(y \mid x)$ according to Eq. \ref{eq:op_policy}. We will argue that there always exists a reward function equivalent to $r$ that admits the desired reparameterization. We construct the projection $f$ as  
\begin{equation}
    f(r; \pi_\text{ref}, \beta)(x, y) = r(x, y) - \beta\log\sum_{y}\pi_\text{ref}(y\mid x)\exp\left(\frac{1}{\beta}r(x, y)\right)
\end{equation}
Here, $f$ renormalizes $r(x, y)$ by subtracting the logarithm of the partition function of $\pi_r$. The normalization term is dependent solely on $x$ (the prefix), which means $f(r; \pi_\text{ref}, \beta)(x, y)$ belongs to the equivalence class of $r(x, y)$. If we substitute $r$ with the right-hand side of Eq.~\ref{eq:main_eq} (which is valid for any reward), we obtain $f(r; \pi_\text{ref}, \beta)(x, y) = \beta \log \frac{\pi_r(y\mid x)}{\pi_\text{ref}(y\mid x)}$. Thus, this projection produces a representative from $r$'s equivalence class with the required parametric form, establishing that the proposed reparameterization is as general as needed for the reward model.
\end{sproof}

An alternative interpretation of Theorem~\ref{thm:main} is that it precisely characterizes which reward function from each equivalence class is chosen by the DPO reparameterization—specifically, it identifies the reward function that satisfies:

\begin{equation}\label{eq:lag_p}
     \sum_{y}\underbrace{\pi_\text{ref}(y\mid x)\exp\left(\frac{1}{\beta}r(x, y)\right)}_{=\pi(y\mid x)\text{, using Thm.~\ref{thm:main} reparam.}} = 1,
\end{equation}

That is, $\pi(y\mid x)$ forms a valid probability distribution (non-negative values that sum to 1). Notably, by referencing Eq.~\ref{eq:op_policy}, Eq.~\ref{eq:lag_p} can be recognized as the partition function for the optimal policy defined by the reward $r(x, y)$. The central idea behind the DPO algorithm is to introduce explicit constraints to the otherwise under-determined Plackett-Luce family (and, more specifically, the Bradley-Terry model) of preference models. By doing so, one maintains the flexibility of the reward models that can be represented, while ensuring that the optimal policy given by Eq.~\ref{eq:op_policy} becomes analytically computable for any prompt $x$.

\subsection{Instability of Actor-Critic Algorithms}

Our framework also allows us to analyze why conventional actor-critic methods, such as PPO, often encounter instability when used for RLHF. Focusing on the RL fine-tuning stage described in Section~\ref{section:prelims}, we can relate our approach to the control as inference framework introduced in \cite{levine2018reinforcement}, specifically for the constrained RL problem described in \ref{eq:RL}. Assuming we work with a parameterized policy $\pi_{\theta}(y\mid x)$, the aim is to minimize $\mathbb{D}_{\text{KL}}[\pi_{\theta}(y|x) \mid \mid \pi^*(y\mid x)]$, with $\pi^*$ being the optimal policy from Eq.~\ref{eq:optimum_model} associated with reward $r_{\phi}(y, x)$. Through straightforward algebraic manipulation, this leads to the following optimization problem:

\begin{equation}\label{eq:AC}
    \max_{\pi_{\theta}}\mathbb{E}_{\pi_{\theta}(y\mid x)}\bigg[\underbrace{r_{\phi}(x, y) -\beta\log\sum_{y}\pi_\text{ref}(y\mid x)\exp\left(\frac{1}{\beta}r_{\phi}(x, y)\right)}_{f(r_{\phi}, \pi_\text{ref}, \beta)} - \underbrace{\beta\log\frac{\pi_{\theta}(y\mid x)}{\pi_\text{ref}(y\mid x)}}_{\text{KL}}\bigg]
\end{equation}

The objective under consideration aligns with those targeted in earlier studies 
\citep{ziegler2020finetuning, stiennon2022learning, bai2022training, ouyang2022training}, where a reward belonging to the $r_{\phi}$ class is employed, corresponding to a DPO-equivalent framework. Within this context, the normalization component in $f(r_{\phi}, \pi_\text{ref}, \beta)$ can be viewed as the soft value function associated with the reference policy $\pi_\text{ref}$. Although this normalization does not influence the location of the optimum, omitting it could result in a policy gradient with considerable variance, thereby impeding stable learning. Introducing a value function estimator for the normalization can mitigate this, yet such estimation often proves challenging to optimize. Alternatively, earlier methods have addressed normalization by employing a baseline derived from human completions, effectively leveraging a single-sample Monte-Carlo estimate of the normalization factor. In contrast, the DPO reparameterization introduces a reward structure that eliminates the need for such baselines.

\section{Experiments}
Here, we provide empirical analyses of DPO's effectiveness in training policies using only preference data. We begin by investigating a controlled text-generation environment to answer: How does DPO balance the objectives of maximizing reward and minimizing KL-divergence from the reference policy in comparison to prevalent preference optimization techniques like PPO? Following this, we assess DPO's scalability and robustness on more substantial models and challenging RLHF benchmarks, such as summarization and dialogue tasks. Our experiments reveal that, with minimal hyperparameter adjustment, DPO consistently matches or exceeds the performance of established baselines, including RLHF via PPO, as well as selecting the best out of $N$ sampled outputs based on a trained reward function. The forthcoming subsections elaborate on our experimental design; supplementary information is provided in Appendix~\ref{app:exp_details}.

\textbf{Tasks.} We examine three distinct open-ended text generation tasks in our experiments. Across all setups, the learning algorithms derive a policy using a dataset of preferences $\mathcal{D}=\bigl\{x^{(i)}, y_w^{(i)}, y_l^{(i)}\bigr\}_{i=1}^N$. In the case of \textbf{controlled sentiment generation}, $x$ corresponds to a partial movie review from the IMDb corpus \cite{maas-EtAl:2011:ACL-HLT2011}, with the policy tasked with generating a continuation $y$ that reflects positive sentiment. To enable controlled evaluation, we automatically create pairs of preferences for the generations by employing a pre-trained sentiment classifier, ensuring $p(\text{positive}\mid x,y_w)>p(\text{positive}\mid x,y_l)$. For the SFT baseline, GPT-2-large is fine-tuned on the training portion of the IMDB dataset until convergence (see further elaboration in App~\ref{app:sentiment_details}). Within the \textbf{summarization} task, $x$ denotes a Reddit forum post, and the objective is for the policy to produce a summary $y$ highlighting the main arguments. Building on previous approaches, we utilize the Reddit TL;DR summarization dataset \citep{volske-etal-2017-tl} in combination with the human preference annotations obtained by \citeauthor{stiennon2022learning}. The SFT model is fine-tuned on human-constructed summaries of forum threads\footnote{\url{https://huggingface.co/CarperAI/openai_summarize_tldr_sft}} by leveraging the TRLX framework \citep{leandro_von_werra_2023_7790115} for RLHF. The dataset containing human preferences was collected by \citeauthor{stiennon2022learning} using samples produced by another, comparably-trained, SFT model. In the \textbf{single-turn dialogue} scenario, $x$ represents a human-generated query, which might range from scientific inquiries to requests for personal guidance. Here, the policy must deliver a reply $y$ that is both informative and engaging. For this, we employ the Anthropic Helpful and Harmless dialogue dataset \citep{bai2022training}, which includes 170k conversational exchanges between people and automated assistants. Each dialogue concludes with two assistant-generated responses produced by a large (but unspecified) language model, alongside an annotation indicating the response preferred by humans. Unlike previous tasks, a pre-trained SFT model is not present in this context; as a result, we train an SFT model by fine-tuning a standard language model solely on responses identified as preferred.

\textbf{Evaluation.} Our experimental design incorporates two distinct evaluation strategies. To systematically assess how well each algorithm addresses the constrained reward maximization objective, we examine performance in the controlled sentiment generation scenario by plotting the set of achievable (reward, KL-divergence) trade-offs. This analysis is possible due to access to the true reward signal, provided by a sentiment classifier. In contrast, practical applications often lack an explicit ground-truth reward. Therefore, in real-world-inspired summarization and single-turn dialogue tasks, we quantify algorithmic performance in terms of their \textit{win rate} against a baseline policy. To facilitate this, we rely on GPT-4 as a stand-in for human raters, employing it to evaluate summary quality and helpfulness of responses, respectively. The baselines consist of reference summaries from the test corpus in the summarization case, and the dataset's preferred responses for dialogue. While prior work has shown that large language models may outperform traditional automated metrics as evaluators \citep{Chen2023ExploringTU}, we independently validate our reliance on GPT-4 through a dedicated human study (see Sec.~\ref{sec:human-judgments}). Results demonstrate that GPT-4's judgments exhibit strong alignment with human assessments, with observed human-GPT-4 agreement often matching or exceeding the level of consistency between human annotators.

\textbf{Methods.} Beyond DPO, our evaluation encompasses a suite of baseline and comparison methods for aligning language model outputs with human preferences. For summarization, we test zero-shot prompting with \textbf{GPT-J} \citep{gpt-j}, whereas in the dialogue domain, we use 2-shot prompting of \textbf{Pythia-2.8B} \citep{biderman2023pythia}. We further benchmark a standard \textbf{SFT} model alongside \textbf{Preferred-FT}, which is trained using supervised learning on preferred completions $y_w$. The sources for $y_w$ differ by task—coming from the SFT model in both sentiment-controlled and summarization experiments, or from a general-purpose LM in dialogue. Another supervised-inspired approach evaluated is \textbf{Unlikelihood}~\citep{welleck2019neural}, which drives the policy to boost the likelihood of $y_w$ while explicitly penalizing the model’s probability for $y_l$; an adjustable coefficient $\alpha\in[0,1]$ is applied to this penalty component. We also analyze \textbf{PPO} \citep{schulman2017proximal}, which employs a reward model fitted from preference annotations, as well as \textbf{PPO-GT}, an oracle version leveraging the true reward in sentiment-controlled settings. In the latter case, two PPO-GT variants are compared: a standard off-the-shelf implementation \cite{leandro_von_werra_2023_7790115}, and an enhanced version featuring reward normalization and fine-tuned hyperparameters—these enhancements are also applied to standard PPO using learned rewards. Lastly, we include the \textbf{Best of $N$} baseline: for each input, we sample $N$ outputs from the SFT model (or from Preferred-FT for dialogue), selecting the response that scores highest under a preference-trained reward model. While this approach isolates reward model quality from PPO optimization and can achieve strong results, it remains computationally costly, requiring $N$ model generations per test prompt even for moderate values of $N$.

\subsection{How well can DPO optimize the RLHF objective?}

In standard RLHF methods, the objective of maximizing the KL-constrained reward seeks a compromise between maximizing reward and minimizing divergence from a reference policy. Consequently, comparisons between algorithms should consider not only the attained reward but also the KL divergence; improvements in reward accompanied by disproportionately higher KL values may not be advantageous. Figure~\ref{fig:frontier-tldr-main} illustrates the tradeoff frontier between reward and KL for a selection of algorithms evaluated on the sentiment task. For each algorithm, we conduct several training sessions, each with a distinct parameter value controlling the degree of policy conservatism (target KL values of $\{3,6,9,12\}$ for PPO, $\beta$ settings of $\{0.05,0.1,1,5\}$, and $\alpha$ options of $\{0.05,0.1,0.5,1\}$ for unlikelihood; for preferred-FT, we vary random seeds). This comprehensive sweep comprises 22 separate runs. Following every 100 training iterations up to convergence, each trained policy is evaluated on a held-out set of prompts, where we calculate both the mean reward using the true reward metric and the average sequence-level KL\footnote{Specifically, the sum of KL divergences computed at each timestep in the sequence.} with respect to the reference policy $\text{KL}\left(\pi\mid \mid \pi_\text{ref}\right)$. The experiments show that DPO establishes a clearly superior frontier, yielding the greatest rewards at comparably low KL distances. This outcome is noteworthy for several reasons. First, even though DPO and PPO are optimizing an identical objective, DPO demonstrates a marked advantage in efficiency, with its reward/KL tradeoff always outperforming PPO. Second, DPO outperforms PPO even when PPO is provided with access to the exact reward values (PPO-GT).

\subsection{Can DPO scale to real preference datasets?}
\label{sec:dpo-real-datasets}
We proceed to assess the effectiveness of DPO for fine-tuning on tasks involving summarization and single-turn dialogue. It is well-documented that standard automatic metrics like ROUGE often show weak alignment with human judgment in summarization tasks~\citep{stiennon2022learning}. Previous research has demonstrated that language models fine-tuned with PPO to align with human preferences generate summaries rated as more helpful. To benchmark various approaches, we sample outputs from the test portion of the TL;DR summarization dataset and determine the mean win rate relative to reference summaries in the test set. For each method, generations are sampled at a range of temperatures between 0.0 and 1.0, and corresponding win rates are presented in Figure~\ref{fig:frontier-tldr-main} (right). DPO, PPO, and Preferred-FT all employ the identical GPT-J SFT checkpoint as their initialization\footnote{\url{https://huggingface.co/CarperAI/openai_summarize_tldr_sft}}. Our analysis shows that DPO achieves a win rate near 61\% when the temperature is set to 0.0, outperforming PPO, which attains about 57\% at its best (also at temperature 0.0). Additionally, DPO secures a greater peak win rate than the best-of-$N$ reference method. It is important to mention that $\beta$, a key DPO hyperparameter, was not extensively optimized; as such, the observed results might not fully represent DPO’s upper bound. Furthermore, DPO’s win rate demonstrates considerably less sensitivity to temperature than that of PPO, whose performance can deteriorate to the level of the unmodified GPT-J model at higher temperatures. The Preferred-FT approach does not offer substantial improvements over the baseline SFT model. A direct human preference evaluation between DPO and PPO (see Section~\ref{sec:human-judgments}) shows that at temperature 0.25, DPO outputs were favored 58\% of the time over PPO outputs at temperature 0.

For the evaluation on single-turn dialogue, we focus on the portion of the Anthropic HH dataset \citep{bai2022training} test split that involves a single round of interaction between human and assistant. To assess different approaches, GPT-4 is utilized to determine win rates, using the test set’s preferred completions as the gold-standard reference. Given the absence of a standard SFT baseline for this scenario, we initialize with a pre-trained Pythia-2.8B model, then employ Preferred-FT to fit a reference model to the curated completions so that the outputs remain in-distribution, followed by DPO training. For further comparison, we benchmark against the strongest result obtained from 128 Preferred-FT completions (the Best of $N$ baseline, where performance saturates at $N=128$ for this dataset, as illustrated in Appendix Figure~\ref{fig:best-of-n}), and also a 2-shot prompted variant of the Pythia-2.8B base model; our findings indicate DPO achieves comparable or superior outcomes for each method’s optimal temperature. We additionally test an RLHF model—trained with PPO on the Anthropic HH dataset \footnote{\url{https://huggingface.co/reciprocate/ppo_hh_pythia-6B}} and released by a reputable group \footnote{\url{https://github.com/CarperAI/trlx/tree/main/examples/hh}}—but cannot identify any prompt or temperature settings that outperform the Pythia-2.8B base model. Given our observations on the TL;DR task and the shared objective optimized by both methods, we treat Best of 128 as an approximate surrogate for PPO model performance. Notably, DPO stands out as the only approach that is both computationally tractable and able to surpass the preferred completions baseline on the Anthropic HH benchmark, delivering results on par with or superior to the resource-intensive Best of 128 strategy. As depicted in Figure~\ref{fig:dialogue-main}, DPO also achieves peak performance after relatively few training steps.

\subsection{Generalization to a new input distribution}

To further assess how PPO and DPO perform under distributional changes, we evaluate the policies obtained from the Reddit TL;DR summarization experiments on a novel input set—specifically, news articles found in the test portion of the CNN/DailyMail dataset \citep{nallapati-etal-2016-abstractive}. The optimal sampling temperatures identified on TL;DR (0 and 0.25) are applied for this evaluation. The outcomes are detailed in Table~\ref{tab:ood}. Using the identical GPT-4 (C) evaluation prompt as in the Reddit TL;DR setting—with the modification of ``forum post'' to ``news article''—we calculated GPT-4's win rate relative to dataset ground-truth summaries. On this distribution, the DPO policy maintains a substantial performance advantage over PPO. These results provide preliminary support that DPO-trained policies generalize at least as effectively as those trained with PPO, despite DPO not leveraging the additional set of unlabeled Reddit TL;DR prompts incorporated by PPO.

\subsection{Validating GPT-4 judgments with human judgments}
\label{sec:human-judgments}
To assess the trustworthiness of GPT-4-based evaluations, we conducted a user study referencing results from the TL;DR summarization task and testing two distinct GPT-4 prompts. The \textbf{GPT-4 (S)} (simple) prompt asks only which summary captures the key points most effectively. By contrast, the \textbf{GPT-4 (C)} (concise) prompt queries which summary is both better and more concise; we include this prompt after noting that GPT-4 (under the simple prompt) prefers verbose, repetitive summaries more than human raters do. The full text of both prompts can be found in Appendix~\ref{app:prompts}. Our evaluation involved three pairwise comparisons representing high (DPO, temp. 0.25), low (PPO, temp. 1.0), and intermediate (SFT, temp. 0.25) quality, aiming to sample from across the quality spectrum; each method was compared to greedily-sampled PPO at its best-performing temperature. Results show that, for both prompts, GPT-4 aligns with human preferences at rates similar to the agreement among human raters themselves. This implies that GPT-4 serves as a credible surrogate for human assessment (with multiple human annotations collected for the DPO and PPO-1 pairs only due to the small rater pool). Notably, the \textbf{GPT-4 (C)} prompt tends to yield win rates more consistent with those from humans; accordingly, we employ this prompt for the core results in Section~\ref{sec:dpo-real-datasets}. For comprehensive information on the human evaluation process—including the web tool shown to raters and the list of participants—refer to Appendix~\ref{app:human-study}.

\section{Discussion}
Preference-based learning presents an effective and adaptable approach for developing language models that are both competent and aligned with human values. In this work, we proposed DPO, a straightforward methodology that leverages preference data for model training without relying on reinforcement learning paradigms. Instead of fitting the preference-learning challenge into a traditional RL mold to utilize pre-existing RL algorithms, DPO establishes a direct correspondence between policy functions and reward representations for language models. This allows models to be optimized according to human preferences using a basic cross-entropy objective, thus eliminating the necessity for reinforcement learning and avoiding any sacrifice in expressivity. With minimal need for hyperparameter adjustment, DPO achieves performance on par with or surpassing current RLHF methods, including those that use PPO. As a result, DPO substantially lowers the complexity of fine-tuning language models with human feedback.

\textbf{Limitations \& Future Work.} Several avenues remain open for exploration based on our findings. It is important to investigate how DPO-trained policies behave on inputs that differ from those seen during training, especially in comparison to models trained using explicit reward functions. Preliminary results indicate comparable generalization abilities between DPO and PPO-based approaches, but thorough investigations are required to draw robust conclusions. For instance, the potential of self-labeling strategies using the DPO policy to effectively leverage data with missing labels warrants further study. Additionally, examining how reward over-optimization appears within the direct preference optimization context is necessary, particularly considering the modest performance decline observed in Figure~\ref{fig:dialogue-main}-right—does this represent such an effect? While our experiments focus on models up to 6B parameters, scaling DPO to substantially larger, state-of-the-art architectures remains a promising research direction. Concerning evaluations, our analysis reveals that win rates produced by GPT-4 are sensitive to prompt variations; further research could focus on optimizing automated evaluation methodologies. Lastly, DPO’s applications likely extend well beyond language model preference alignment and may prove valuable for training generative models across diverse domains.